% answers_short.tex —— 2005 年·普通高等学校招生考试·全国卷Ⅲ·文科 · 简短答案速查
% 由 tools/make_short_answers.py 自 answers.tex 的 \gksolhead/\gkanswer 提取，
% 并与 2005/2005_数学真题_全国卷Ⅲ文科_答案要点.md 逐题比对后生成，勿手改（改源 MD 后重跑）。

\begin{gknotice}
\textbf{使用说明}：本节只给最终答案便于快速核对，完整推导见解析版。答案为编写者独立推导，\textbf{非官方标准答案}。
\end{gknotice}

\gkpart{选择题}
\begin{center}\begin{tabular}{ccccccccccccc}
\toprule
题号 & 1 & 2 & 3 & 4 & 5 & 6 & 7 & 8 & 9 & 10 & 11 & 12 \\
\midrule
答案 & D & B & B & C & A & C & C & B & C & D & D & A \\
\bottomrule
\end{tabular}\end{center}

\gkpart{填空题}
\begin{itemize}[leftmargin=2.6em]
  \item[\textbf{13.}] $3$。
  \item[\textbf{14.}] $k = -\dfrac{2}{3}$。
  \item[\textbf{15.}] $x + y - 2 = 0$。
  \item[\textbf{16.}] $3$。
\end{itemize}

\gkpart{解答题}
\begin{itemize}[leftmargin=2.6em]
  \item[\textbf{17.}] $\left(0, \dfrac{3\pi}{4}\right) \cup \left(\pi, \dfrac{7\pi}{4}\right)$。
  \item[\textbf{18.}] (1) $P(\text{\text{甲}}) = 0.2$，$P(\text{\text{乙}}) = 0.25$，$P(\text{\text{丙}}) = 0.5$；(2) $0.7$
  \item[\textbf{19.}] (1) 由平面 $VAD \perp$ 底面及 $AB \perp AD$ 得 $AB \perp$ 平面 $VAD$；(2) $\arccos\dfrac{\sqrt{21}}{7}$
  \item[\textbf{20.}] $k_n = 3^{n+1}$。（源面印作“等差中项”，对 $d\neq0$ 自相矛盾，解析按“等比中项”处理，见解析正文与 \texttt{missing-issues.md}。）
  \item[\textbf{21.}] 高为 $10\ \mathrm{cm}$ 时容积最大，最大容积 $19600\ \mathrm{cm}^3$。
  \item[\textbf{22.}] (1) 当且仅当 $x_1 + x_2 = 0$ 时直线 $l$ 经过焦点 $F$；(2) $x - 4y + 41 = 0$
\end{itemize}
